\section{Outcome, scope, and novelty standard}

This report proposes four additional directions. The first three produce
local source reductions; the fourth makes two geometric proofs share one
well-ordering argument, at a small increase in total source size. This
distinction matters: a shorter application site is not automatically a
shorter development once the supporting theorem is counted.

\begin{enumerate}
\item \textbf{Identify the existing list code with Mathlib's encoding.}
Use that identity to obtain length and injectivity facts, and derive typed
parent-list correctness from the raw checker theorem already present.
\item \textbf{Base quadratic coefficient extraction on a local pair fibre.}
Prove the coefficient identity for an arbitrary weight map with exactly the
two required contributing ordered pairs, then specialize to global Sidon and
selective coding.
\item \textbf{Remove sign cases from the fundamental Pell-unit proof.}
Work with the magnitude of the second coordinate, retaining the existing
square-gap comparison. Keep the more general minimum-coordinate
criterion as an optional checked abstraction.
\item \textbf{Factor the least-exit argument in the two strip-wall proofs.}
Separate interval minimality from the asymmetric tile geometry on the left
and right of Start.
\end{enumerate}

\subsection{The source snapshot}

Repository access began by resolving the current main branch through GitHub.
All source analysis and integration patches use the immutable commit
\sourcecommit. Its recorded author time is 26 October 2026, 16:43:16 UTC.
The commit already merges several earlier simplification adoptions, but the
adoption table also records proposals that are pending or deferred
\cite{sourcepin,adoption}. We do not interpret those statuses as permission
to repeat them.

The task's premise is stronger than ``find edits missing from main.'' We
treat the mechanisms proposed by earlier reports as exhausted even where
their code is absent, uncompiled, rejected, or only described as future
work. The numerical before/after measurements are nevertheless against the
real pinned files. They are not estimates for a hypothetical tree formed by
simultaneously applying every mutually alternative historical design.

\subsection{Exclusion audit}

The audit inventories 329 files in the two requested exclusion directories
and 400 members of the five incoming ZIPs. The six placed packets are
\code{2026-10-26a} through \code{2026-10-26f}. All eleven principal
manuscripts were inspected by section; supplementary proposal notes,
patches, and the earlier reports' own novelty audits were also searched and
reviewed. The archive contents were read as evidence, not executed.

The accompanying \code{evidence/exclusion-ledger.json} and
\code{evidence/EXCLUSION_LEDGER.md} identify 60 prior technique families,
with file-and-line or archive-member-and-line references.
\path{evidence/prior-corpus-manifest.json} records their source hashes.
Inventory coverage does not assert that every incidental copied source
snapshot or binary received an independent semantic reading.

Examples of deliberately excluded ideas include scanner simulation,
primitive-recursion congruence, decoder completeness, coefficient envelopes,
sum-of-squares interfaces, generated circuit certificates, support envelopes,
finite invariant predecessor lemmas, carry uniqueness, rational-root descent,
and aligned-cell additive equivalences. Some apparently fresh directions
were explicitly rejected in the older manuscripts themselves; those too
are excluded. The present four proposals were compared with their nearest
prior mechanisms, not merely searched under their proposed names.

\subsection{What is claimed}

The supplied patches preserve the existing executable encoding and arithmetic
definitions. New statements are proved, and existing public theorem types
are retained. The formal checks concern proof simplifications; they do not
establish an additional Diophantine equivalence, improve a variable count,
or resolve a mathematical gap elsewhere in the research program.

The report separates three kinds of evidence:
\begin{itemize}
\item \emph{Full modified module checks}: the actual source module is
compiled with its replacement proof and production imports.
\item \emph{Focused checks}: exact theorem applications, helper checks,
and axiom reports examine the changed interface.
\item \emph{Architectural recommendations}: a checked general theorem may
support future reuse without any claim that a whole downstream migration
has already been implemented.
\end{itemize}

All final counts and check outcomes are recorded in
\code{evidence/results.json}. The validation section states the actual
rebuild scope and the limits of the evidence.
